\bta{2009}

\begin{enumerate}
	%\renewcommand{\labelenumi}{\arabic{enumi}.}
	% A(\Alph) a(\alph) I(\Roman) i(\roman) 1(\arabic)
	%设定全局标号series=example	%引用全局变量resume=example
	%[topsep=-0.3em,parsep=-0.3em,itemsep=-0.3em,partopsep=-0.3em]
	%可使用leftmargin调整列表环境左边的空白长度 [leftmargin=0em]
	\item
	填空题 (本题满分 30 分，每小题 6 分)
	\begin{enumerate}
		%\renewcommand{\labelenumi}{\arabic{enumi}.}
		% A(\Alph) a(\alph) I(\Roman) i(\roman) 1(\arabic)
		%设定全局标号series=example	%引用全局变量resume=example
		%[topsep=-0.3em,parsep=-0.3em,itemsep=-0.3em,partopsep=-0.3em]
		%可使用leftmargin调整列表环境左边的空白长度 [leftmargin=0em]
		\item
		设 $y=\sqrt[4]{x \sqrt[3]{e^{x} \sqrt{\sin \frac{1}{x}}}}$, 则 $\frac{d y}{d x}=$  \tk{$y\left(\frac{1}{4 x}+\frac{1}{12}-\frac{1}{24} \cdot \frac{1}{x^{2}} \cdot \cot \frac{1}{x}\right)$} 。
        \nynote{对数求导法，这里我算的好几次都是$y\left(\frac{1}{4 x}+\frac{1}{12}-\frac{1}{24} \cdot \frac{1}{x^{2}} \cdot \cot \frac{1}{x}\right)$，不知是否答案有误。}
		\item  
设 $z=\frac{1}{x} f(x y)+y g(x+y)$, 其中 $f, g$ 具有二阶连续导数, 则 $\frac{\partial^{2} z}{\partial x^{2}}=$  \tk{$\frac{2}{x^{3}} f(x y)-\frac{2 y}{x^{2}} f^{\prime}(x y)+$ $\frac{y^{2}}{x} f^{\prime \prime}(x y)+y g^{\prime \prime}(x+y)$}  。
		
		
		\item
		函数 $f(x)=\frac{1}{x^{2}-2 x-3}$ 展开为 $x$ 的幂级数, 展开式为: \tk{$-\frac{1}{12} \sum\limits_{n=0}^{\infty}\left(\frac{x}{3}\right)^{n}-\frac{1}{4} \sum\limits_{n=0}^{\infty}(-1)^{n} x^{n}$} 。



\item
过点 $P(-1,0,4)$ 且与平面 $3 x-4 y+z+10=0$ 平行， 又与直线 $L: \frac{x+1}{1}=$ $\frac{y-3}{1}=\frac{z}{2}$ 相交的直线方程为:  \tk{$\frac{x+1}{16}=\frac{y}{19}=\frac{z-4}{28}$}  。
\nynote{
    在L上取一点$Q(x,y,z)$，有$\mathbf{QP}$与平面法向量垂直，可解出$Q$点，即得直线方程，
    这里我算的是$\frac{x+1}{8}=\frac{y}{17}=\frac{z-4}{4}$。
}

\item
微分方程 $y y^{\prime\prime}+\left(y^{\prime}\right)^{2}=0$ 满足初始条件 $\left.y\right|_{x=0}=1,\left.y^{\prime}\right|_{x=0}=\frac{1}{2}$ 的解为:  
\tk{
$y^{2}=x+1$
} 。
		
	
		
	\end{enumerate}


\banswer{
\begin{enumerate}
	%\renewcommand{\labelenumi}{\arabic{enumi}.}
	% A(\Alph) a(\alph) I(\Roman) i(\roman) 1(\arabic)
	\item
	取对数再求导。
	\[
	\ln y=\frac{1}{4}\ln x+\frac{1}{12}x+\frac{1}{24}\ln\sin\frac1x ,
	\]
	两端对 $x$ 求导，注意 $\left(\frac1x\right)'=-\frac1{x^{2}}$：
	\[
	\frac{y'}{y}=\frac{1}{4x}+\frac{1}{12}
	+\frac{1}{24}\cdot\frac{\cos\frac1x}{\sin\frac1x}\cdot\left(-\frac{1}{x^{2}}\right),
	\]
	故
	\[
	\frac{dy}{dx}=y\left(\frac{1}{4x}+\frac{1}{12}
	-\frac{1}{24}\cdot\frac{1}{x^{2}}\cot\frac1x\right).
	\]

	\item
	$z=\frac1xf(xy)+yg(x+y)$ 对 $x$ 求一次导：
	\[
	\frac{\partial z}{\partial x}=-\frac{1}{x^{2}}f(xy)+\frac{y}{x}f'(xy)+yg'(x+y) ;
	\]
	再求一次（乘积法则，链式因子仍为 $y$ 与 $1$）：
	\begin{align*}
	\frac{\partial^{2}z}{\partial x^{2}}
	&=\frac{2}{x^{3}}f(xy)-\frac{y}{x^{2}}f'(xy)-\frac{y}{x^{2}}f'(xy)\\
	&\quad+\frac{y^{2}}{x}f''(xy)+yg''(x+y)\\
	&=\frac{2}{x^{3}}f(xy)-\frac{2y}{x^{2}}f'(xy)
	+\frac{y^{2}}{x}f''(xy)+yg''(x+y).
	\end{align*}

	\item
	先部分分式：$\frac{1}{x^{2}-2x-3}=\frac{1}{(x-3)(x+1)}
	=\frac14\left(\frac{1}{x-3}-\frac{1}{x+1}\right)$。再分别展开
	\[
	\frac{1}{x-3}=-\frac13\cdot\frac{1}{1-\frac x3}=-\frac13\sum_{n=0}^{\infty}\left(\frac x3\right)^{n},
	\qquad
	\frac{1}{x+1}=\sum_{n=0}^{\infty}(-1)^{n}x^{n} .
	\]
	故
	\[
	f(x)=-\frac{1}{12}\sum_{n=0}^{\infty}\left(\frac x3\right)^{n}
	-\frac14\sum_{n=0}^{\infty}(-1)^{n}x^{n},
	\]
	两个级数分别在 $|x|<3$、$|x|<1$ 内收敛，故展开式对 $|x|<1$ 成立。

	\item
	设所求直线与 $L$ 交于 $Q$。把 $L$ 写成参数式 $Q=(-1+t,\ 3+t,\ 2t)$，则
	$\vv{PQ}=(t,\ 3+t,\ 2t-4)$。直线与平面平行等价于 $\vv{PQ}$ 垂直于法向量
	$\bm{n}=(3,-4,1)$：
	\[
	3t-4(3+t)+(2t-4)=t-16=0\quad\Longrightarrow\quad t=16 .
	\]
	于是方向向量 $\vv{PQ}=(16,19,28)$，所求直线为
	\[
	\frac{x+1}{16}=\frac{y}{19}=\frac{z-4}{28} .
	\]
	（校验：$(16,19,28)\cdot(3,-4,1)=48-76+28=0$；又 $P$ 不在平面上
	（$3(-1)-0+4+10=11\neq0$），故确为“平行”而非“含于”。）

	\item
	注意 $\left(yy'\right)'=\left(y'\right)^{2}+yy''=0$，故 $yy'\equiv C$。
	由初值 $y(0)=1$、$y'(0)=\frac12$ 得 $C=\frac12$，即
	$\left(y^{2}\right)'=2yy'=1$，积分得 $y^{2}=x+C'$；
	再由 $y(0)=1$ 得 $C'=1$，故
	\[
	y^{2}=x+1,\qquad y=\sqrt{x+1} .
	\]
\end{enumerate}

	\par\vspace{1mm}\noindent
	\textbf{注：}{\kaishu
	第 (1) 小题原答案作 $\tan\frac1x$，有误：
	$\left(\ln\sin u\right)'=\cot u$，故应为 $\cot\frac1x$。
	以 $x=0.7$ 数值求导得 $\frac{y'}{y}=0.428300$，而含 $\cot$ 的表达式给出 $0.428300$、
	含 $\tan$ 的给出 $-0.153371$，确认原答案有误。

	第 (4) 小题原有“可能有问题，后期重算”的批注，现已复核：答案
	$\frac{x+1}{16}=\frac{y}{19}=\frac{z-4}{28}$ 正确。
	旁注中提出的 $\frac{x+1}{8}=\frac{y}{17}=\frac{z-4}{4}$ 虽与 $L$ 相交，
	但 $(8,17,4)\cdot(3,-4,1)=-40\neq0$，并不平行于所给平面，故不合题意。
	}\par
}


\item
选择题 (本题满分 30 分，每小题 6 分)
\begin{enumerate}
	%\renewcommand{\labelenumi}{\arabic{enumi}.}
	% A(\Alph) a(\alph) I(\Roman) i(\roman) 1(\arabic)
	%设定全局标号series=example	%引用全局变量resume=example
	%[topsep=-0.3em,parsep=-0.3em,itemsep=-0.3em,partopsep=-0.3em]
	%可使用leftmargin调整列表环境左边的空白长度 [leftmargin=0em]
	\item
	$ \lim\limits _{x \rightarrow 2} \frac{\int_{2}^{x}\left(\int_{t}^{2} e^{-u^{2}} du \right) d t}{(x-2)^{2}}=$ \xzanswer{D} 
	
	
\fourchoices
{$\frac{1}{e^{2}}$}
{$-\frac{1}{e^{2}}$}
{$\frac{1}{2 e^{4}}$}
{$-\frac{1}{2 e^{4}}$}

	
	
	\item
	求积分 $\int_{0}^{1} e^{\sqrt{1-x}} d x=$ \xzanswer{C} 
	
	
\fourchoices
{$ 0 $}
{$ 1 $}
{$ 2 $}
{$ 3 $}


\item
设在 $[0,+\infty)$ 上 $f^{\prime \prime}(x)>0$, 则当 $x \in(0,+\infty)$ 时, 下列不等式成立的是 \xzanswer{B} 


\fourchoices
{$f^{\prime}(0) x<f(0)-f(x)<f^{\prime}(x) x$}
{$f^{\prime}(0) x<f(x)-f(0)<f^{\prime}(x)x$}
{$f(0)-f(x)>f^{\prime}(0) x>f^{\prime}(x)x$}
{$f(0)-f(x)<f^{\prime}(0) x<f^{\prime}(x)x$}




\item
设 $L: x^{2}+(y+1)^{2}=2$ 取逆时针方向, 则 $\oint_{L} \frac{x d y-y d x}{x^{2}+(y+1)^{2}}=$ \xzanswer{B} 


\fourchoices
{$4 \pi$}
{$2 \pi$}
{$\pi^{2}$}
{$2 \pi^{2}$}



\item
设 $\Sigma$ 为曲面 $z=x^{2}+y^{2}(0 \leqslant z \leqslant 1)$ 的下侧, 则曲面积分 $I=\iint\limits_{\Sigma} y^{3} d z d x+(y+z) d x d y=$ \xzanswer{C} 


\fourchoices
{$-\frac{\pi}{2}$}
{$\frac{\pi}{2}$}
{$-\frac{\pi}{4}$}
{$\frac{\pi}{4}$}

\nynote{这里我用定义法和补顶面再高斯定理的方法算出来都是A.$-\pi/2$，不知是否答案有误。}

\end{enumerate}


\banswer{
\begin{enumerate}
	%\renewcommand{\labelenumi}{\arabic{enumi}.}
	% A(\Alph) a(\alph) I(\Roman) i(\roman) 1(\arabic)
	\item
	\textbf{D}。分子分母同趋于 $0$，连用两次洛必达法则：
	\[
	\lim_{x\to2}\frac{\int_{x}^{2}e^{-u^{2}}du}{2(x-2)}
	=\lim_{x\to2}\frac{-e^{-x^{2}}}{2}=-\frac{1}{2e^{4}} .
	\]
	（第一次求导用变限积分求导，注意外层被积函数在 $t=x$ 处的值是 $\int_x^2e^{-u^2}du$。）

	\item
	\textbf{C}。令 $u=\sqrt{1-x}$，则 $x=1-u^{2}$，$dx=-2u\,du$，积分限 $x:0\to1$ 对应 $u:1\to0$，
	\[
	\int_{0}^{1}e^{\sqrt{1-x}}dx=2\int_{0}^{1}ue^{u}du
	=2\left[ue^{u}-e^{u}\right]_{0}^{1}=2 .
	\]

	\item
	\textbf{B}。$f''>0$ 说明 $f'$ 严格递增，$f$ 是凸函数。对 $x>0$，
	由拉格朗日中值定理存在 $\xi\in(0,x)$ 使 $f(x)-f(0)=f'(\xi)x$，
	而 $f'(0)<f'(\xi)<f'(x)$，故
	\[
	f'(0)x<f(x)-f(0)<f'(x)x .
	\]

	\item
	\textbf{B}。作平移 $Y=y+1$，则 $L$ 变为 $(x,Y)$ 平面上以原点为心、$\sqrt2$ 为半径的圆，且
	\[
	x\,dy-y\,dx=\left(x\,dY-Y\,dx\right)+dx .
	\]
	第一部分是标准的 $\oint\frac{x\,dY-Y\,dx}{x^{2}+Y^{2}}=2\pi$；
	第二部分 $\oint\frac{dx}{x^{2}+Y^{2}}=\frac12\oint dx=0$（沿闭曲线且分母恒为 $2$）。
	故原积分 $=2\pi$。

	\item
	\textbf{C}。$\Sigma:\ z=x^{2}+y^{2}$ 取下侧，$d\bm{S}=\left(z_x,z_y,-1\right)dxdy$，
	即 $dydz\to2x\,dxdy$、$dzdx\to2y\,dxdy$、$dxdy\to-dxdy$。投影区域 $D:\ x^{2}+y^{2}\leqslant1$，
	\begin{align*}
	I&=\iint\limits_{D}\left[2y^{4}-\left(y+x^{2}+y^{2}\right)\right]dxdy\\
	&=\frac{\pi}{4}-0-\frac{\pi}{2}=-\frac{\pi}{4} ,
	\end{align*}
	其中用了 $\iint_Dy^{4}=\frac{\pi}{8}$、$\iint_Dy=0$、$\iint_D\left(x^{2}+y^{2}\right)=\frac{\pi}{2}$。

	也可补顶面 $\Sigma_{1}:\ z=1,\ x^{2}+y^{2}\leqslant1$ 取上侧用高斯公式：
	$\operatorname{div}=3y^{2}+1$，
	\[
	\iiint_{\Omega}\left(3y^{2}+1\right)dV=\frac{\pi}{4}+\frac{\pi}{2}=\frac{3\pi}{4},
	\]
	而 $\iint_{\Sigma_{1}}=\iint_{D}(y+1)dxdy=\pi$（顶面上 $dzdx=0$），
	故 $I=\frac{3\pi}{4}-\pi=-\frac{\pi}{4}$。
\end{enumerate}

	\par\vspace{1mm}\noindent
	\textbf{注：}{\kaishu
	第 (5) 小题旁注称两种方法都算得 $-\frac{\pi}{2}$（选 A），实为计算疏漏：
	上面两种方法均给出 $-\frac{\pi}{4}$，二重积分数值计算也给出 $-0.785398=-\frac{\pi}{4}$。
	常见的错误是漏掉 $2y^{4}$ 这一项（它来自 $y^{3}dzdx$），
	或在高斯公式中漏掉 $3y^2$ 对应的 $\frac{\pi}{4}$。原答案 C 正确。
	}\par
}


\item
(本题满分 10 分)
设 $\lim\limits _{x \rightarrow \infty}\left(\frac{x-a}{x+a}\right)^{x}=\int_{a}^{+\infty} x e^{-2 x} d x$, 求 $a$ 的值。



\banswer{
	\textbf{左端.}
	\[
	\left(\frac{x-a}{x+a}\right)^{x}=\left(1-\frac{2a}{x+a}\right)^{x}
	=\left[\left(1-\frac{2a}{x+a}\right)^{-\frac{x+a}{2a}}\right]^{-\frac{2ax}{x+a}} ,
	\]
	当 $x\to\infty$ 时中括号内 $\to e$，指数 $\to-2a$，故左端极限为 $e^{-2a}$。

	\textbf{右端.} 分部积分（$a$ 为任意实数，积分收敛）：
	\begin{align*}
	\int_{a}^{+\infty}\!xe^{-2x}dx
	&=\left[-\frac x2e^{-2x}\right]_{a}^{+\infty}+\frac12\int_{a}^{+\infty}\!e^{-2x}dx\\
	&=\frac a2e^{-2a}+\frac14e^{-2a}=\frac{2a+1}{4}e^{-2a} .
	\end{align*}

	\textbf{解方程.} 两端相等，约去 $e^{-2a}>0$：
	\[
	\frac{2a+1}{4}=1\quad\Longrightarrow\quad a=\frac{3}{2} .
	\]
}


\item 
(本题满分 10 分)
设 $f(x)$ 在 $[0,1]$ 上连续, $f(x)>0$, 求证：存在唯一的 $a \in(0,1)$, 使得 $\int_{0}^{a} f(t) d t=\int_{a}^{1} \frac{1}{f(t)} d t$。

\banswer{
令 $F(x)=\int_0^x f(t)dt-\int_x^1\frac1{f(t)}dt$，$x\in[0,1]$。
$F$ 连续，且 $F(0)=-\int_0^1\frac1{f(t)}dt<0$，$F(1)=\int_0^1 f(t)dt>0$。
又 $F'(x)=f(x)+\frac1{f(x)}>0$（$x\in(0,1)$），故 $F$ 严格单调递增。
由连续函数的介值定理，存在唯一的 $a\in(0,1)$ 使得 $F(a)=0$，
即 $\int_0^a f(t)dt=\int_a^1\frac1{f(t)}dt$。
}



\item 
(本题满分 10 分)
由直线 $x=0, y=8$ 及抛物线 $y=x^{2}$ 围成一个曲边三角形, 在曲边 $y=x^{2}$ 上求一点 $M(X, Y)$, 使得在该点处的切线与直线 $x=0, y=8$ 所围成的三角形面积最小。

\banswer{
	由 $y=x^{2}$ 关于 $y$ 轴的对称性，只需讨论 $a>0$ 一侧（对应曲边三角形
	$0\leqslant x\leqslant2\sqrt2$），另一侧取相反数即可。

	\textbf{写出三角形的面积.} 设切点 $M\left(a,a^{2}\right)$，切线
	\[
	y=2a(x-a)+a^{2}=2ax-a^{2} .
	\]
	它与 $x=0$ 交于 $\left(0,-a^{2}\right)$，与 $y=8$ 交于
	$\left(\frac{8+a^{2}}{2a},8\right)$；再与 $x=0,\ y=8$ 的交点 $(0,8)$ 一起
	构成一个直角三角形，两直角边长分别为 $8+a^{2}$ 与 $\frac{8+a^{2}}{2a}$。故
	\[
	S(a)=\frac{\left(8+a^{2}\right)^{2}}{4a}
	=\frac{16}{a}+4a+\frac{a^{3}}{4} .
	\]

	\textbf{求极小.}
	\[
	S'(a)=-\frac{16}{a^{2}}+4+\frac{3a^{2}}{4}=0
	\ \Longrightarrow\ 3a^{4}+16a^{2}-64=0 ,
	\]
	记 $w=a^{2}>0$ 得 $3w^{2}+16w-64=0$，正根 $w=\frac83$，即
	\[
	a=\sqrt{\frac83}=\frac{2\sqrt6}{3}\approx1.633\in\left(0,2\sqrt2\right).
	\]
	$S'$ 在此处由负变正，故为极小值点，也是 $(0,+\infty)$ 上的最小值点。

	\textbf{结论.} $M$ 的坐标为
	$\left(\frac{2\sqrt6}{3},\frac83\right)$ 或 $\left(-\frac{2\sqrt6}{3},\frac83\right)$，
	\[
	S_{\min}=\frac{\left(8+\frac83\right)^{2}}{4\cdot\frac{2\sqrt6}{3}}
	=\frac{64\sqrt6}{9}\approx17.42 .
	\]
}



\item 
(本题满分 10 分)
求满足 $x=\int_{0}^{x} f(t) d t + \int_{0}^{x} t f(t-x) d t$ 的可微函数 $f(x)$ 。

\nynote{该题在2007年重复出现}
\banswer{
	\textbf{化去卷积.} 在 $\int_{0}^{x}tf(t-x)dt$ 中令 $u=t-x$，得
	\[
	\int_{0}^{x}\!tf(t-x)dt=\int_{-x}^{0}\!(u+x)f(u)du
	=\int_{-x}^{0}\!uf(u)du+x\int_{-x}^{0}\!f(u)du .
	\]
	于是原方程成为
	\[
	x=\int_{0}^{x}\!f(t)dt+\int_{-x}^{0}\!uf(u)du+x\int_{-x}^{0}\!f(u)du .
	\]

	\textbf{第一次求导.} 注意 $\frac{d}{dx}\int_{-x}^{0}g(u)du=g(-x)$，故
	\[
	1=f(x)-xf(-x)+\int_{-x}^{0}\!f(u)du+xf(-x) ,
	\]
	含 $xf(-x)$ 的两项相消，得关系式
	\[
	f(x)+\int_{-x}^{0}\!f(u)du=1 .
	\]
	令 $x=0$ 得 $f(0)=1$。

	\textbf{第二次求导.} 对上式求导得
	\[
	f'(x)+f(-x)=0 ,
	\]
	令 $x=0$ 得 $f'(0)=-f(0)=-1$。再对上式求导一次：
	\[
	f''(x)-f'(-x)=0 ,
	\]
	而在 $f'(x)=-f(-x)$ 中以 $-x$ 代 $x$ 得 $f'(-x)=-f(x)$，代入即
	\[
	f''(x)+f(x)=0 .
	\]

	\textbf{解方程.} 通解 $f=A\cos x+B\sin x$；由 $f(0)=1$、$f'(0)=-1$ 得 $A=1,\ B=-1$，
	\[
	f(x)=\cos x-\sin x .
	\]
	（代回原方程，两端之差恒为 $0$，已用符号计算验证。）

	\par\vspace{1mm}\noindent
	\textbf{注：}{\kaishu
	关键一步是从 $f'(x)+f(-x)=0$ 里消去 $f(-x)$。除上面的“以 $-x$ 代 $x$”外，
	也可直接由关系式 $f(x)+\int_{-x}^{0}f=1$ 解出 $\int_{-x}^{0}f$ 再求导，效果相同。
	本题与 2007 年真题重复。
	}\par
}


\item 
(本题满分 10 分)
证明方程 $F\left(x+\frac{z}{y}, y+\frac{z}{x}\right)=0$ 所确定的隐函数 $z=z(x, y)$ 满足方程 $x \frac{\partial z}{\partial x}+y \frac{\partial z}{\partial y}=z-x y_{\circ}$。


\banswer{
	记 $u=x+\dfrac{z}{y}$，$v=y+\dfrac{z}{x}$，则 $F(u,v)=0$。
	两端分别对 $x$、$y$ 求偏导（$z=z(x,y)$）：
	\begin{align*}
	F_u\left(1+\frac{z_x}{y}\right)+F_v\left(\frac{z_x}{x}-\frac{z}{x^{2}}\right)&=0,\\
	F_u\left(\frac{z_y}{y}-\frac{z}{y^{2}}\right)+F_v\left(1+\frac{z_y}{x}\right)&=0 .
	\end{align*}

	\textbf{解出偏导.} 记 $A=\dfrac{F_u}{y}+\dfrac{F_v}{x}=\dfrac{xF_u+yF_v}{xy}$（设 $A\neq0$），
	整理上两式得
	\[
	z_xA=\frac{zF_v}{x^{2}}-F_u,\qquad
	z_yA=\frac{zF_u}{y^{2}}-F_v .
	\]

	\textbf{组合.} 于是
	\begin{align*}
	x\frac{\partial z}{\partial x}+y\frac{\partial z}{\partial y}
	&=\frac{z\left(\frac{F_v}{x}+\frac{F_u}{y}\right)-\left(xF_u+yF_v\right)}{A}\\
	&=z-\frac{xF_u+yF_v}{A}=z-xy ,
	\end{align*}
	最后一步用了 $\dfrac{xF_u+yF_v}{A}=xy$。证毕。

	\par\vspace{1mm}\noindent
	\textbf{注：}{\kaishu
	推导中要求 $xy\neq0$ 且 $A=\frac{xF_u+yF_v}{xy}\neq0$，
	后者正是隐函数定理成立（$z$ 可解出且可微）的条件。
	本题与 2023 年选择题第 (8) 小题实质相同。
	}\par
}


\newpage
\item 
(本题满分 10 分)
已知函数 $y_{1}=x e^{x}+e^{2 x}, y_{2}=x e^{x}+e^{-x}, y_{3}=x e^{x}+e^{2 x}-e^{-x}$ 是某二阶常系数线性非齐次微分方程的三个解，求此微分方程的表达式。



\banswer{
	\textbf{提取齐次解.} 非齐次方程两解之差是齐次方程的解：
	\[
	y_{1}-y_{3}=e^{-x},\qquad y_{1}-y_{2}=e^{2x}-e^{-x} ,
	\]
	故 $e^{-x}$ 与 $e^{2x}$ 都是齐次方程的解，且线性无关。

	\textbf{定齐次方程.} 特征根为 $-1$ 与 $2$，特征方程
	$(r+1)(r-2)=r^{2}-r-2=0$，即齐次方程为 $y''-y'-2y=0$。

	\textbf{定右端.} 三个解的公共部分 $y^{*}=xe^{x}$ 是非齐次方程的一个特解。
	由 $y^{*\prime}=(1+x)e^{x}$、$y^{*\prime\prime}=(2+x)e^{x}$ 得
	\begin{align*}
	y^{*\prime\prime}-y^{*\prime}-2y^{*}
	&=\left[(2+x)-(1+x)-2x\right]e^{x}\\
	&=(1-2x)e^{x} .
	\end{align*}

	\textbf{结论.}
	\[
	y''-y'-2y=(1-2x)e^{x} .
	\]
	（把 $y_1,y_2,y_3$ 逐一代入，残差均为 $0$，已用符号计算验证。）
}


\item 
(本题满分 10 分)
设 $f(x)$ 是周期为 2 的周期函数，且 $f(x)=\left\{\begin{array}{ll}x & 0 \leqslant x \leqslant 1 \\ 0 & 1<x<2\end{array}\right.$, 试求出 $f(x)$ 的傅里叶级数,
并求数项级数 $\sum\limits_{n=0}^{\infty} \frac{1}{(2 n+1)^{2}}$ 的和。

\nynote{若代入1求解，则要注意$x=1$是一个跳跃间断点，傅里叶展开式在此处的值应该是$\frac{1}{2}$。}
\banswer{
	\textbf{先写出一个周期.} 周期为 $2$，取 $(-1,1)$ 作基本区间。
	当 $-1<x<0$ 时 $x+2\in(1,2)$，故 $f(x)=f(x+2)=0$；当 $0\leqslant x\leqslant1$ 时 $f(x)=x$。
	于是半周期 $l=1$，
	\[
	f(x)\sim\frac{a_{0}}{2}+\sum_{n=1}^{\infty}\left(a_{n}\cos n\pi x+b_{n}\sin n\pi x\right).
	\]

	\textbf{算系数.} 由 $f$ 在 $(-1,0)$ 上为零，所有积分只剩 $\int_0^1$：
	\[
	a_{0}=\int_{-1}^{1}\!f\,dx=\int_{0}^{1}\!x\,dx=\frac12,\qquad \frac{a_0}{2}=\frac14 ;
	\]
	分部积分得
	\begin{align*}
	a_{n}&=\int_{0}^{1}\!x\cos n\pi x\,dx=\frac{(-1)^{n}-1}{n^{2}\pi^{2}},\\
	b_{n}&=\int_{0}^{1}\!x\sin n\pi x\,dx=\frac{(-1)^{n+1}}{n\pi} .
	\end{align*}
	故
	\[
	f(x)=\frac14+\sum_{n=1}^{\infty}\frac{(-1)^{n}-1}{n^{2}\pi^{2}}\cos n\pi x
	+\sum_{n=1}^{\infty}\frac{(-1)^{n+1}}{n\pi}\sin n\pi x .
	\]
	该式在 $f$ 的连续点处成立，在跳跃点处收敛到左右极限的平均值。

	\textbf{求级数和.} $x=0$ 是连续点（左右极限都是 $0$，且 $f(0)=0$），代入：
	\[
	0=\frac14+\sum_{n=1}^{\infty}\frac{(-1)^{n}-1}{n^{2}\pi^{2}} .
	\]
	求和中只有奇数 $n$ 有贡献，且贡献为 $-\frac{2}{n^{2}\pi^{2}}$，故
	\[
	\frac14=\frac{2}{\pi^{2}}\sum_{n=0}^{\infty}\frac{1}{(2n+1)^{2}}
	\ \Longrightarrow\
	\sum_{n=0}^{\infty}\frac{1}{(2n+1)^{2}}=\frac{\pi^{2}}{8} .
	\]
	（数值：$\sum_{n<2\times10^{6}}=1.2337004$，$\frac{\pi^{2}}{8}=1.2337006$。）

	\par\vspace{1mm}\noindent
	\textbf{注：}{\kaishu
	若改用 $x=1$ 代入，须注意 $x=1$ 是跳跃间断点（$f(1^{-})=1$，$f(1^{+})=0$），
	级数在该点收敛到 $\frac{1+0}{2}=\frac12$ 而不是 $f(1)=1$；
	这样得到 $\frac12=\frac14+\frac{2}{\pi^{2}}\sum\frac{1}{(2n+1)^{2}}$，结论相同。
	}\par
}



\item 
(本题满分 10 分)
设 $f(x)$ 连续, 令 $F(t)=\iint\limits_{x^{2}+y^{2} \leqslant t^{2}} f\left(x^{2}+y^{2}\right) d x d y$ $(t \geqslant 0)$ 。 求 $F^{\prime \prime}(0)$ 。

\banswer{
	\textbf{化为单积分.} 用极坐标，
	\[
	F(t)=\int_{0}^{2\pi}\!\!d\theta\int_{0}^{t}\!f\left(r^{2}\right)r\,dr
	=2\pi\int_{0}^{t}\!f\left(r^{2}\right)r\,dr
	=\pi\int_{0}^{t^{2}}\!f(u)du ,
	\]
	最后一步作了代换 $u=r^{2}$。

	\textbf{一阶导.} 由 $f$ 连续及变限积分求导，
	\[
	F'(t)=\pi f\left(t^{2}\right)\cdot2t=2\pi tf\left(t^{2}\right),\qquad F'(0)=0 .
	\]

	\textbf{二阶导.} $f$ 只假设连续，不能再对 $F'$ 求导，须回到定义（$t\geqslant0$，故为右导数）：
	\[
	F''(0)=\lim_{t\to0^{+}}\frac{F'(t)-F'(0)}{t}
	=\lim_{t\to0^{+}}\frac{2\pi tf\left(t^{2}\right)}{t}=2\pi f(0) .
	\]

	\par\vspace{1mm}\noindent
	\textbf{注：}{\kaishu
	这里不可以直接写 $F''(t)=2\pi f\left(t^2\right)+4\pi t^{2}f'\left(t^{2}\right)$，
	因为题设只给了 $f$ 连续，$f'$ 未必存在；用导数定义才是稳妥的做法。
	}\par
}


\item 
(本题满分 10 分)
设 $f(x)$ 在闭区间 $[0,1]$ 上连续， 在开区间 $(0,1)$ 内可导，且 $f(0)=0$，$ f(1)=1$ 。试证明 任意给定的正数 $a, b$ 在开区间 $(0,1)$ 内存在不同的实数 $\xi$ 和 $\eta$, 使得 
\[ 
\frac{a}{f^{\prime}(\xi)}+\frac{b}{f^{\prime}(\eta)}=a+b
 \]

\banswer{
	\textbf{取分点.} 因 $a,b>0$，有 $0<\dfrac{a}{a+b}<1$，即
	$f(0)=0<\dfrac{a}{a+b}<1=f(1)$。由 $f$ 在 $[0,1]$ 上连续及介值定理，存在 $c\in(0,1)$ 使
	\[
	f(c)=\frac{a}{a+b} .
	\]

	\textbf{两次拉格朗日中值定理.} 分别在 $[0,c]$ 与 $[c,1]$ 上用中值定理：存在
	$\xi\in(0,c)$、$\eta\in(c,1)$ 使
	\begin{align*}
	f'(\xi)&=\frac{f(c)-f(0)}{c-0}=\frac{a}{(a+b)c},\\
	f'(\eta)&=\frac{f(1)-f(c)}{1-c}=\frac{b}{(a+b)(1-c)} .
	\end{align*}
	两者都是正数，特别地不为零，倒数有意义。

	\textbf{合成.}
	\begin{align*}
	\frac{a}{f'(\xi)}+\frac{b}{f'(\eta)}
	&=a\cdot\frac{(a+b)c}{a}+b\cdot\frac{(a+b)(1-c)}{b}\\
	&=(a+b)c+(a+b)(1-c)=a+b .
	\end{align*}
	又 $\xi<c<\eta$，故 $\xi\neq\eta$，命题得证。
}
	






	
\end{enumerate}


